\documentclass[12 pt]{article}

\usepackage{stoversymb}
\usepackage[margin=0.75in, top=0.875in, bottom=0.875in]{geometry}
\usepackage{amsmath,amsfonts,amssymb,url,multicol,graphicx,tikz,soul,changepage}
%===makes urls render well===
\usepackage{lmodern}
\usepackage[T1]{fontenc}
%============================
\usepackage{wasysym} % smileys
\usepackage[inline]{enumitem}

\everymath{\displaystyle}

%\setenumerate{itemsep=0.25in}
\setlist[enumerate,1]{label=\arabic*., itemsep=0.5in}
\setlist[enumerate,2]{label={(\alph*)}, itemsep=0.5in}
\setlist[enumerate,3]{label={\roman*.}}
\setlist[itemize,1]{label=$\circ$, itemsep=0.25in}

\newcommand{\truefalse}[1]{#1\hfill\rule[-1mm]{220pt}{0.75pt}}
\newcommand{\hint}[1]{\vspace{3mm}\textbf{Hint}: #1}
\newcommand{\note}[1]{\textbf{Note}: #1}
\newcommand{\infsum}[3]{\sum_{{#1}={#2}}^\infty {#3}}
\newcommand{\axes}[1]{\begin{center}\includegraphics[scale=#1]{3DAxes}\end{center}}
\newcommand{\comps}[1]{\langle #1_1,#1_2,#1_3\rangle}
\newcommand{\compslong}[3]{\langle #1, #2, #3\rangle}
\newcommand{\ijk}[2]{#1\vect{#2}}
\newcommand{\pmat}[1]{\begin{pmatrix}#1\end{pmatrix}}
\newcommand{\bmat}[1]{\begin{bmatrix}#1\end{bmatrix}}

\graphicspath{ {./../img/} }
\DeclareGraphicsExtensions{.pdf}

\begin{document}
\begin{flushright}Name: \line(1,0){200}\end{flushright}
\begin{center}
\Large{\textbf{MAS 3105 --- Homework 1}}
\end{center}
\textbf{Directions:} Complete the following problems for a homework grade. Solutions \textit{must} be presented in a neat and professional manner in order to receive credit, answers given without showing work will not be eligible to receive partial credit, and \textit{work for the problems \textbf{must} be done on scratch paper and not on this handout!} \textbf{Date Due:} Thursday, February 1.
\vspace{0.125in}
\begin{enumerate}[leftmargin=0in, rightmargin=-0.25in]
%	%========Slack==============
	\item \note{Yes, you will get graded for this question. \smiley}
	\begin{enumerate}[itemsep=0.25in]
		\item Navigate to our course homepage at
		\begin{center} \url{http://www.math.fsu.edu/~cstover/teaching/sp18_mas3105/}
		\end{center}
		\item Read and familiarize yourself with the three resources listed under \textit{Supplementary Resources} on the \textsc{General Info} tab.
		\item Follow the instructions for using \textsc{Slack} messenger.\begin{quote}\textbf{Note:} This may require that I approve your email address, so to avoid some last minute glitch where I don't get to your approval on-time, please don't wait to do this!\end{quote}
		\item Navigate to the channel \url{#introductions} in the left column under \textsc{Channels}; its browser URL should be something like \url{https://spring2018-linalg.slack.com/messages/C8WQS8E77/}.
		\item Introduce yourself by \ul{answering each of the following questions}:
		\begin{quote}\textbf{Note:} This will be visible to everyone who signs into our class's chat room, so you definitely want to keep these answers PG-13, safe for work, and non-incriminatory. {\Large\smiley}\end{quote}
		\begin{enumerate}[label=(\roman*),itemsep=3mm]
			\item What is your name?
			\item Where are you from (using any interpretation you'd like)?
			\item What is the best (interpret this however you'd like) place you've ever visited/lived? Why is it so special to you?
			\item How long have you been in Tallahassee?
			\item What is your major?
			\item What do you like to do for fun? (besides linear algebra, of course!)
			\item What is the coolest math/science ``thing'' you know? Why is it interesting to you?
		\end{enumerate}
		\item Under which username did you register for \textsc{Slack}?\hspace{6mm}\line(1,0){200}
	\end{enumerate}
	
	\item Put each of the following matrices in reduced row echelon form (RREF). 
	
	\textbf{Note:} If you get to row echelon form (REF) \textit{before} you get to RREF, please indicate the REF matrix row equivalent to the one you started with!
	\begin{enumerate}
		\item $\sansA=\pmat{2 & 3 \\ 1 & 5}$
		\item $\sansB=\pmat{0 & 2 & 3 \\ 0 & 1 & 5 \\ 0 & 0 & -1}$
		\item $\sansC=\pmat{-3 & 6 & 9 \\ 1 & -2 & -3 \\ 0 & 5 & 1 \\ 0 & 0 & -8}$
		\item $\sansD=\pmat{2 & 1 & 0 & -1 \\ -1 & 1 & 1 & 3 \\ 0 & 1 & -3 & 5 \\ 0 & 0 & 1 & 1}$
	\end{enumerate}
	
	\item For each of the following linear systems, decide whether it's (i) consistent with a unique-solution, (ii) consistent with a non-unique solution, or (iii) inconsistent. If a system is consistent, solve it using matrices; otherwise, state why it isn't and/or how you know it's not.
	\begin{adjustwidth}{0in}{-0.375in}
		\begin{multicols}{3}
			\begin{enumerate}[itemsep=0.75in,leftmargin=0in,rightmargin=0in]
				\item $\begin{array}{ccccc}x_1 & + & x_2 & = & 3 \\ 2x_1 & - & x_2 & = & 1\end{array}$
				\item $\begin{array}{ccccccc}x_1 & + & & & x_3 & = & 3 \\ 2x_1 & - & x_2 & & & = & 1\end{array}$
				\item $\begin{array}{ccccccc}x_1 & + & x_2 & + & x_3 & = & 1 \\ -x_1 & - & x_2 & + & x_3 & = & 1 \\ x_1 & & & &  & = & 2\end{array}$
				\item $\begin{array}{ccccccc}x_1 & + & x_2 & + & x_3 & = & 1 \\ -x_1 & - & x_2 & + & x_3 & = & 1 \\ & & & &  x_3 & = & 0\end{array}$
				\item $\begin{array}{ccccccc}x_1 & + & x_2 & + & x_3 & = & 1 \\ -x_1 & - & x_2 & + & x_3 & = & -1 \\ & & & &  x_3 & = & 0\end{array}$
				\item $\begin{array}{ccccccc}3x_1 & - & x_2 & + & 2x_3 & = & 7 \\ 4x_1 & - & 4x_2 & + & x_3 & = & -1 \\ x_1 & - & 3x_2 & + & x_3 & = & 1\end{array}$
			\end{enumerate}
		\end{multicols}
	\end{adjustwidth}

	\newpage
	
	\item For each of the six systems of equations in exercise 3, write a vector equation that is equivalent.\vspace{1in}
	
	\item For each of the six systems of equations in exercise 3, write a matrix equation that is equivalent.\vspace{1in}

	\item Each of the following matrices corresponds to the \textbf{augmented matrix} of a linear system. Determine whether the corresponding system is (i) consistent with a unique-solution, (ii) consistent with a non-unique solution, or (iii) inconsistent. \textbf{Justify your claim!}
	\begin{adjustwidth}{0in}{-0.375in}
		\begin{multicols}{2}
			\begin{enumerate}[itemsep=1in]
				\item $\sansA=\pmat{1 & 1 & 3 \\ 0 & 1 & -1}$
				\item $\sansB=\pmat{0 & 2 & -3 & 1 \\ 1 & 1 & 1 & 1 \\ 0 & 0 & 0 & 1}$
				\item $\sansC=\pmat{0 & 2 & -3 & 1 \\ 1 & 1 & 1 & 1 \\ 0 & 0 & 0 & 0}$
				\item $\sansD=\pmat{0 & 2 & -3 & 1 \\ 1 & 1 & 1 & 1 \\ -2 & 1 & 3 & 11 \\ 0 & 0 & 0 & 0}$\vspace{0.5in}
				\item $\sansE=\pmat{1 & 2 & 3 & 4 \\ 5 & 6 & 7 & -8 \\ -9 & 10 & -11 & 12}$
			\end{enumerate}
		\end{multicols}
	\end{adjustwidth}

	\newpage
	
	\item Is $\vect{b}=\pmat{2 \\ -1 \\ 6}$ a linear combination of the vectors $\vect{a}_1=\pmat{1 \\ 1 \\ 0}$, $\vect{a}_2=\pmat{0 \\ 1 \\ -1}$, and $\vect{a}_3=\pmat{2 \\ 0 \\ -1}$? Why or why not?\vspace{0.375in}
	
	\item Is $\vect{b}=\pmat{-4 \\ -2 \\ 1}$ a linear combination of the vectors formed by the columns of \mbox{$\sansA=\pmat{-1 & 1 & 3 \\ 1 & 1 & 2 \\ 0 & 1 & -4}$}? Why or why not?\vspace{0.375in}
	
	\item \textbf{True or False:} Every vector in $\Reals^2$ is a linear combination of the vectors $\vect{u}=\pmat{1 \\ 0}$ and $\vect{v}=\pmat{0 \\ 2}$. \textbf{Justify your claim!}
	
	\hint{The clearest way to do this is to start with an arbitrary vector $\vect{w}=\pmat{a \\ b}$ in $\Reals^2$ and to either write $\vect{w}=s\vect{u}+t\vect{v}$ for real numbers $s$ and $t$ \ul{or} to show that no such expression is possible.}\vspace{0.375in}
	
	\item Describe all solutions of $\sansA\vect{x}=\vect{0}$ in parametric vector form, where $\sansA$ is row equivalent to the given matrix.
	\begin{enumerate}
		\item $\pmat{1 & 3 & -3 & 4 \\ 0 & 1 & -4 & 5}$
		\item $\pmat{1 & 0 & 0 \\ 0 & 1 & 0 \\ 0 & 0 & 34}$
		\item $\pmat{1 & 5 & 2 & -6 & 9 & 0 \\ 0 & 0 & 1 & -7 & 4 & -8 \\ 0 & 0 & 0 & 0 & 0 & 1 \\ 0 & 0 & 0 & 0 & 0 & 0}$
	\end{enumerate}	
	
	\item Let 
	\[\sansA=\pmat{1 & 0 & -2 \\ 0 & 3 & 6 \\ 4 & 2 & -3}\quad\quad\text{and}\quad\quad\vect{b}=\pmat{4 \\ -1 \\ 4},\]
	write $\vect{a}_1$, $\vect{a}_2$, and $\vect{a}_3$ for the columns of $\sansA$, and let $W=\vsspan\{\vect{a}_1,\vect{a}_2,\vect{a}_3\}$.
	\begin{enumerate}
		\item Write the augmented matrix for the linear system that corresponds to the matrix equation $\sansA\vect{x}=\vect{b}$.
		\item Solve the matrix equation from (a) and write the solution as a vector.
		\item Is $\vect{b}$ in $\{\vect{a}_1,\vect{a}_2,\vect{a}_3\}$?
		\item How many vectors are in $\{\vect{a}_1,\vect{a}_2,\vect{a}_3\}$?
		\item Is $\vect{b}$ in $W$?
		\item How many vectors are in $W$?
		\item Show that $\vect{a}_2$ is in $W$.
		\item Show that $\vect{a}_1-\vect{a}_3$ is in $W$.
		\item Do the columns of $\sansA$ form a linearly independent set? Why or why not?
		\item Let $h$ be a real number and let $\vect{u}=\pmat{1 \\ -1 \\ h}$. For which value(s) $h$ is $\vect{u}$ in $\vsspan\{\vect{a}_1,\vect{a}_2\}$?
		\item Let $\vect{u}$ be as in (j). For what value(s) of $h$ is $\{\vect{a}_2,\vect{b},\vect{u}\}$ linearly \textit{dependent}?
	\end{enumerate}
\end{enumerate}
\end{document}